\themaD
\graphicspath{{../../S06_Interpreter_representer_des_donnees/Images/}}



\chapter{Interpréter,\\représenter\\des données}
\label{S06}

%%%%%%%%%%%%%%%%%%%%%%%%%%%%%%%%%%
%%%%%%%%%%%%%%%%%%%%%%%%%%%%%%%%%%
\begin{autoeval}
   \small
   \begin{enumerate}
      \item Il recueille et organise des données.
      \item Il lit et interprète des données brutes ou présentées sous forme de tableaux, de diagrammes et de graphiques.
      \item Il représente, sur papier ou à l’aide d’un tableur-grapheur, des données sous la forme d’un tableau, d’un diagramme ou d’un graphique.
   \end{enumerate}
\end{autoeval}

\begin{prerequis}
   \begin{itemize}
      \item[\com] Recueillir des données, les organiser.
      \item[\com] Lire et interpréter des données sous forme de données brutes, de tableau, de diagramme (diagramme en bâtons, diagramme circulaire, histogramme).
      \item[\com] Utiliser un tableur-grapheur pour présenter des données sous la forme d’un tableau ou d’un diagramme.
   \end{itemize}
\end{prerequis}

\vfill


\begin{debat}[Débat : le premier tableur]
   Des données brutes récoltées ont souvent peu de sens si elles sont utilisées ainsi, d'où la nécessité de les disposer d'une manière plus lisible à l'aide de tableaux et diagrammes. \\
   Avec l'avènement de l'informatique, les tableaux deviennent numériques grâce à l'apport des {\bf tableurs} : logiciels qui permettent de manipuler des données numériques, d'effectuer un certain nombre d'opérations de façon automatisée, de créer des représentations graphiques à partir des données : diagrammes , histogrammes, courbes\dots{} \\
   Le premier tableur fut créé en 1978 par {\it Daniel Bricklin}, étudiant à Harvard qui devait établir des tableaux comptables pour une étude de cas sur Pepsi-Cola sans pour autant établir tous les calculs \og à la main \fg. Son premier prototype, {\it VisiCalc} (pour Visible Calculator), pouvait manipuler un tableau de vingt lignes et cinq colonnes ! \\
   \begin{center}
      \begin{Tableur}[Bandeau=false,Colonnes=5]
         & & & & \\
         & & & & \\
         & & & & \\
      \end{Tableur}
   \end{center}
   \bigskip
   \begin{cadre}[B2][J4]
      \begin{center}
         Vidéo : \href{https://www.ted.com/talks/dan_bricklin_meet_the_inventor_of_the_electronic_spreadsheet#t-556741}{\bf Meet the inventor of electronic spreadsheet}, site Internet de {\it Ted talks}.
      \end{center}
   \end{cadre}
\end{debat}


%%%%%%%%%%%%%%%%%%%%%%%%%%%%%%%%
%%%%%%%%%%%%%%%%%%%%%%%%%%%%%%%%
\activites

\begin{activite}[Récolter des données]
   {\bf Objectif :} Récolter des données au sein de la classe afin de les représenter sous différentes formes. \\
   \begin{QCM}
      De combien d'enfants est composée ta famille proche (éventuellement recomposée), toi y compris ?
      \begin{center}
      {\hautab{2.5}
      \begin{ttableau}{0.9\linewidth}{4}
         \hline
         & & & \\
         \hline
         & & & \\
         \hline
         & & & \\
         \hline
         & & & \\
         \hline
         & & & \\
         \hline
         & & & \\
         \hline
         & & & \\
         \hline
      \end{ttableau}}
      \end{center}
      \ \\
      Ces données seront utilisées pour construire la trace écrite.
   \end{QCM}
\end{activite}


%%%%%%%%%%%%%%%%%%%%%%%
%%%%%%%%%%%%%%%%%%%%%%%
\cours 


%%%%%%%%%
\section{Tableaux}

On souhaite connaître le nombre d'enfants qui constitue une fratrie d'une classe de 5\up{e} du collège Simone Veil composée ce jour là de 25 élèves. On obtient les résultats suivants :
   \begin{center}
      1 ; 2 ; 5 ; 1 ; 2 ; 2 ; 3 ; 4 ; 1 ; 3 ; 2 ; 6 ; 2 ; 3 ; 4 ; 2 ; 6 ; 1 ; 3 ; 2 ; 1 ; 2 ; 3 ; 3 ; 2.
   \end{center}
   
On {\bf organise} les résultats : pour rassembler les données, on les présente sous forme d'un tableau où l'on regroupe ensemble les différentes valeurs obtenues. L'{\bf effectif} d'une donnée est le nombre de fois qu'elle apparaît.
\begin{center}
   \Stat[Sondage,Tableau,Largeur=8mm,Donnee=Nombre d'enfants,Total]{1,2,5,1,2,2,3,4,1,3,2,6,2,3,4,2,6,1,3,2,1,2,3,3,2}
\end{center}


%%%%%%%%%%%
\section{Diagrammes en barres}

\begin{definition}
   Un \textbf{diagramme en barres} (ou en bâtons) est un graphique qui représente une série par de barres rectangulaires avec des hauteurs ou des longueurs proportionnelles aux effectifs.
\end{definition}

\bigskip

\begin{minipage}{8cm}
   {\it Diagramme en barres verticales}\Stat[Graphique,Grille,Unitex=0.9,Unitey=0.4,EpaisseurBatons=8,Donnee=Nombre d'enfants,ListeCouleursB={LightSeaGreen,0.15[LightSeaGreen,white],0.3[LightSeaGreen,white],0.45[LightSeaGreen,white],0.6[LightSeaGreen,white],0.75[LightSeaGreen,white]}]{1/5,2/9,3/6,4/2,5/1,6/2}
\end{minipage}
\begin{minipage}{8cm}
   {\it Diagramme en barres horizontales}
      \Stat[Qualitatif,Graphique,Barre,Grille,Longueur=6cm,EcartBarre=2mm,ListeCouleurs={LightSeaGreen,0.15[LightSeaGreen,white],0.3[LightSeaGreen,white],0.45[LightSeaGreen,white],0.6[LightSeaGreen,white],0.75[LightSeaGreen,white]}]{1 enfant/5,2 enfants/9,3 enfants/6,4 enfants/2,5 enfants/1,6 enfants/2}
\end{minipage}


\begin{definition}
   Un \textbf{diagramme circulaire} est un graphique qui représente une série par des secteurs circulaires dont la mesure des angles est proportionnelle aux effectifs.
\end{definition}
   
\bigskip
   
\begin{minipage}{8.6cm}
   Pour déterminer la valeur d'un angle, on effectue un calcul de proportionnalité : \\
   Les 25 familles sont représentées sur le diagramme par un secteur angulaire de \udeg{360}. \\
   Une famille correspond à un secteur égal à $\dfrac{\udeg{360}}{25} =\udeg{14,4}$. \\
   Il suffit donc de multiplier l'effectif par 14,4 pour obtenir la valeur de l'angle (ici arrondie au degré). \\ [2mm]
   \Stat[Qualitatif,Tableau,Angle,Donnee=Nombre d'enfants,Largeur=0.5cm]{1/5,2/9,3/6,4/2,5/1,6/2}
\end{minipage}
\qquad
\begin{minipage}{7cm}
   {\it Diagramme circulaire} \\ [5mm]
   \Stat[Qualitatif,Graphique,Angle,Rayon=1.8cm,ListeCouleurs={LightSeaGreen,0.15[LightSeaGreen,white],0.3[LightSeaGreen,white],0.45[LightSeaGreen,white],0.6[LightSeaGreen,white],0.75[LightSeaGreen,white]}]{1 enfant/5,2 enfants/9,3 enfants/6,4 enfants/2,5 enfants/1,6 enfants/2}
\end{minipage}


%%%%%%%%%%%%%%%%%%%%%%%%%%%%%%%%%%%
%%%%%%%%%%%%%%%%%%%%%%%%%%%%%%%%%%%
\exercicesbase

\begin{colonne*exercice}

\begin{exercice} %1
   On a demandé à des élèves de trois classes de 5\up{e} combien d'animaux de compagnie vivaient avec eux. \\
   Le résultat est représenté dans le diagramme suivant :
   \Stat[Graphique,Grille,Unitex=1.2,Unitey=0.14,Origine=-0.5,Pasy=5,EpaisseurBatons=8,Donnee=,Effectif=Nombre d'élèves,LectureFine,Tiret,ListeCouleursB={Gray,Gray,Gray,Gray,Gray}]{0/32,1/21,2/15,3/4,4/3} \\ [-8mm]
   \hspace*{27mm} Nombre d'animaux de compagnie \\ [2mm]
   Les affirmations suivantes sont-elles vraies ou fausses ?
   \begin{enumerate}
      \item 21 élèves ont un seul animal de compagnie.
      \item Il y a 75 élèves en 5\up{e} dans ce collège.
      \item Les élèves qui ont deux animaux de compagnie sont trois fois plus nombreux que les élèves qui en ont trois.
      \item 70 élèves ont moins de trois animaux de compagnie.
      \item Plus de la moitié des élèves ont au moins un animal de compagnie.
      \item Parmi les élèves qui ont au moins un animal de compagnie, la moitié en ont plusieurs.
   \end{enumerate}
\end{exercice}

\begin{corrige}
\ \\ [-5mm]
   \begin{enumerate}
      \item {\blue vrai}, par lecture de la deuxième barre.
      \item {\blue faux}. Le diagramme ne nous permet pas de connaître le nombre d'élèves de ce collège puisqu'il s'agit des données de trois classes seulement.
      \item {\blue faux}. 15 élèves ont deux animaux de compagnie et  4 en ont trois. or, 15 n'est pas le triple de 4.
      \item {\blue faux}. $32+21+15 =68$. 68 élèves ont moins de trois animaux de compagnie.   
      \item {\blue vrai}. 32 élèves n'ont pas d'animal de compagnie et 43 en ont au moins un.
      \item {\blue vrai}. Parmi les 43 élèves qui ont au moins un animal de compagnie, 22 en ont plusieurs, donc plus de la moitié.
   \end{enumerate}
\end{corrige}

\bigskip


\begin{exercice} %2
   L'infographie suivante donne le nombre de fois que des sortilèges de magie apparaissent dans les sept livres de la série {\it Harry Potter}.
   \begin{center}
      \includegraphics[width=8cm]{infographie}
   \end{center}
   {\footnotesize\it Source de l'infographie : Harry Potter, les nombres d'or de la saga, Le Figaro.fr, 2017}
   \begin{enumerate}
      \item Comment sont représentées les données dans cette infographie ?
      \item Combien y a-t-il eu de sortilèges donnés au total dans les sept livres ?
      \item Représenter les données dans un tableau en notant uniquement les sortilèges cités plus de 20 fois.
      \item Construire un diagramme en bâtons pour les valeurs de ce tableau.
      \item Construire un diagramme circulaire pour les valeurs de ce tableau.
      \item Laquelle de ces trois représentations vous convient le mieux ? Pourquoi ?
   \end{enumerate}
\end{exercice}

\begin{corrige}
   \ \\ [-5mm]
   \begin{enumerate}
      \item Les données sont représentées sous forme de {\blue bulles proportionnelles à l'effectif}.
      \item $140+49+39+34+27+27+24+22+18+15+13+11+11+10+9+7+6+6+4+3 =475$. \\
      {\blue 475 sorts ont été donnés durant les sept livres}.
      \item Tableau des effectifs : \\ \smallskip
         {\small
         \hautab{1}
         \setlength{\tabcolsep}{0cm}
         \begin{Ltableau}{\linewidth}{9}{c}
            \hline       
            Nom & E.P. & Imp. & End. & Accio &Stup. & Exp. & A.K. & Lum. \\
            \hline
            Eff. & 140 & 49 & 39 & 34 & 27 & 27 & 24 & 22 \\
            \hline
         \end{Ltableau}}
      \item Diagramme en bâtons : \\
         {\psset{yunit=0.4,xunit=0.8}
         \begin{pspicture}(-0.8,-1)(8,8.25)
         {\footnotesize
            \psline(0,0)(8,0)
            \multido{\i=0+1,\n=0+20}{8}{\psline(-0.1,\i)(0.1,\i) \rput(-0.5,\i){\n}}
            \psline{->}(0,0)(0,8)
            \psset{linewidth=0.5,linecolor=blue}
            \psline(0.5,0)(0.5,7)
            \psline(1.5,0)(1.5,2.45)
            \psline(2.5,0)(2.5,1.95)
            \psline(3.5,0)(3.5,1.7)
            \psline(4.5,0)(4.5,1.35)
            \psline(5.5,0)(5.5,1.35)
            \psline(6.5,0)(6.5,0.6)
            \psline(7.5,0)(7.5,0.55)
            \rput(0.5,-0.5){E.P.}
            \rput(1.5,-0.6){Imp.}
            \rput(2.5,-0.5){End.}
            \rput(3.5,-0.5){Accio}
            \rput(4.5,-0.6){Stup.}
            \rput(5.5,-0.6){Exp.}
            \rput(6.5,-0.5){A.K.}
            \rput(7.5,-0.5){Lum.}
            \rput(1,7.7){\it Effectifs}}
         \end{pspicture}}
      \item Diagramme circulaire : \\
         {\psset{unit=0.6}
         \footnotesize
         \begin{pspicture}(-5,-3.25)(3,3.25)
            \pscircle(0,0){3}
            \pswedge[fillstyle=solid,fillcolor=blue!70](0,0){3}{0}{139}
            \pswedge[fillstyle=solid,fillcolor=blue!65](0,0){3}{139}{188}
            \pswedge[fillstyle=solid,fillcolor=blue!60](0,0){3}{188}{227}      
            \pswedge[fillstyle=solid,fillcolor=blue!55](0,0){3}{227}{261}
            \pswedge[fillstyle=solid,fillcolor=blue!50](0,0){3}{261}{287}
            \pswedge[fillstyle=solid,fillcolor=blue!45](0,0){3}{287}{314}
            \pswedge[fillstyle=solid,fillcolor=blue!40](0,0){3}{314}{338}
            \pswedge[fillstyle=solid,fillcolor=blue!35](0,0){3}{338}{360}
            \rput(1.5;69){\white Expecto}
            \rput(2;69){\white Patronum}
            \rput{-17}(2;163){\white Impero}
            \rput{28}(2;208){\white Endoloris}
            \rput{62}(2.2;242){\white Accio}
            \rput{274}(2.2;274){\white Stupéfix}
            \rput{300}(1.9;300){\white Expelliarmus}
            \rput{326}(2.2;326){\white A.K.}
            \rput{349}(2.2;349){\white Lumos}
         \end{pspicture}}
      \item Cette question est personnelle\dots
   \end{enumerate}
\end{corrige}

\bigskip


\begin{exercice} %3
   Le tableau ci-dessous représente la répartition des médailles françaises aux Jeux olympiques d'été de 1896 à 2016 pour les dix sports ayant eu le plus de médailles.
   \begin{enumerate}
      \item Compléter le tableau.
      \item Comment est établi le classement des sports aux Jeux olympiques.
      \item Construire trois diagrammes circulaires : celui du cyclisme, du tir et du canoë-kayak en fonction de la couleur de médaille obtenue. Les comparer.
   \end{enumerate}
   \begin{center}
      \small
         {\hautab{2.1}
         \begin{tabular}{|*{7}{c|}}
            \hline 
            Pl. & & Sport & \pscircle[fillstyle=solid,fillcolor=Gold](0,0.1){0.2} & \pscircle[fillstyle=solid,fillcolor=lightgray](0,0.1){0.2} & \pscircle[fillstyle=solid,fillcolor=brown](0,0.1){0.2} & T. \\
            \hline
            1 & \includegraphics[height=4mm]{S1} & & & \, 51 \, & \, 35 \, & \, {\bf 118} \, \\
            \hline 
            2 & \includegraphics[height=4mm]{S2} & & 41 & & 23 & {\bf 91} \\
            \hline
            3 & \includegraphics[height=4mm]{S3} & & 14 & 25 & & {\bf 68} \\
            \hline  
            4 & \includegraphics[height=4mm]{S4} & & 14 & 13 & 10 & \\
            \hline  
            5 & \includegraphics[height=4mm]{S5} & & 14 & 10 & & {\bf 49} \\
            \hline  
            6 & \includegraphics[height=4mm]{S6} & & 13 & & 17 & {\bf 41} \\
            \hline  
            7 & \includegraphics[height=4mm]{S7} & & & 14 & 10 & {\bf 33} \\
            \hline  
            8 & \includegraphics[height=4mm]{S8} & & 9 & & 3 & {\bf 15} \\
            \hline  
            9 & \includegraphics[height=4mm]{S9} & & 8 & 15 & & {\bf 43} \\
            \hline  
            10 & \includegraphics[height=4mm]{S10} & {\textcolor{white}{Canoé-kayak}} & 8 & 9 & 19 & \\
            \hline    
         \end{tabular}} \\ [2mm]
      \hfill {\footnotesize\it Source : France aux Jeux olympiques, Wikipedia, 2019}
   \end{center}
\end{exercice}

\begin{corrige}
   \ \\ [-5mm]
   \begin{enumerate}
      \item Voici le tableau complété : \\ \smallskip
         {\footnotesize
         \hautab{2}
         \begin{tabular}{|*{7}{c|}}
            \hline 
             & & Sport & \pscircle[fillstyle=solid,fillcolor=Gold](0,0.1){0.2} & \pscircle[fillstyle=solid,fillcolor=lightgray](0,0.1){0.2} & \pscircle[fillstyle=solid,fillcolor=brown](0,0.1){0.2} & T. \\
            \hline
            1 & \includegraphics[scale=0.3]{../../S06_Interpreter_representer_des_donnees/Images/S1} & {\blue Escrime} & \, {\blue 32} \, & \, 51 \, & \, 35 \, & \, {\bf 118} \, \\
            \hline 
            2 & \includegraphics[scale=0.3]{../../S06_Interpreter_representer_des_donnees/Images/S2} & {\blue Cyclisme} & 41 & {\blue 27} & 23 & {\bf 91} \\
            \hline
            3 & \includegraphics[scale=0.3]{../../S06_Interpreter_representer_des_donnees/Images/S3} & {\blue Athlétisme} & 14 & 25 & {\blue 29} & {\bf 68} \\
            \hline  
            4 & \includegraphics[scale=0.3]{../../S06_Interpreter_representer_des_donnees/Images/S4} & {\blue Équitation} & 14 & 13 & 10 & {\blue \bf 37} \\
            \hline  
            5 & \includegraphics[scale=0.3]{../../S06_Interpreter_representer_des_donnees/Images/S5} & {\blue Judo} & 14 & 10 & {\blue 25} & {\bf 49} \\
            \hline  
            6 & \includegraphics[scale=0.3]{../../S06_Interpreter_representer_des_donnees/Images/S6} & {\blue Voile} & 13 & {\blue 11} & 17 & {\bf 41} \\
            \hline  
            7 & \includegraphics[scale=0.3]{../../S06_Interpreter_representer_des_donnees/Images/S7} & {\blue Tir} & {\blue 9} & 14 & 10 & {\bf 33} \\
            \hline  
            8 & \includegraphics[scale=0.3]{../../S06_Interpreter_representer_des_donnees/Images/S8} & {\blue Haltérophilie} & 9 & {\blue 3} & 3 & {\bf 15} \\
            \hline  
            9 & \includegraphics[scale=0.3]{../../S06_Interpreter_representer_des_donnees/Images/S9} & {\blue Natation} & 8 & 15 & {\blue 20} & {\bf 43} \\
            \hline  
            10 & \includegraphics[scale=0.3]{../../S06_Interpreter_representer_des_donnees/Images/S10} & {\blue Canoé-kayak} & 8 & 9 & 19 & {\blue \bf 36} \\
            \hline    
         \end{tabular}} \medskip
      \item Le classement des sports est établi grâce au {\blue nombre de médailles d'or}, puis d'argent, puis de bronze.
      \item On récapitule dans un tableau les angles : \\ \smallskip
   \end{enumerate}
      \small
      {\hautab{1.5}
      \begin{Ltableau}{0.95\linewidth}{5}{c}
         \hline   
         Couleur & Or & Argent & Bronze & Total \\
         \hline   
         Cyclisme & 41 & 27 & 23 & 91 \\
         Angle & \textcolor{blue}{\udeg{162}} & \textcolor{blue}{\udeg{107}} & \textcolor{blue}{\udeg{91}} & \udeg{360} \\
         \hline
         Tir & 9 & 14 & 10 & 33 \\
         Angle & \textcolor{blue}{\udeg{98}} & \textcolor{blue}{\udeg{153}} & \textcolor{blue}{\udeg{109}} & \udeg{360} \\
         \hline
         Canoé-kayak & 8 & 9 & 19 & 36 \\
         Angle & \textcolor{blue}{\udeg{80}} & \textcolor{blue}{\udeg{90}} & \textcolor{blue}{\udeg{190}} & \udeg{360} \\
        \hline
      \end{Ltableau}}
      {\psset{unit=0.4}
      \begin{pspicture}(-2,-3.3)(3,3.8)
         \pscircle(0,0){3}
         \pswedge[fillstyle=solid,fillcolor=Gold](0,0){3}{0}{162}
         \pswedge[fillstyle=solid,fillcolor=lightgray](0,0){3}{162}{269}
         \pswedge[fillstyle=solid,fillcolor=brown](0,0){3}{269}{360}      
      \end{pspicture}
      \begin{pspicture}(-3.5,-3)(3,3.5)
         \pscircle(0,0){3}
         \pswedge[fillstyle=solid,fillcolor=Gold](0,0){3}{0}{98}
         \pswedge[fillstyle=solid,fillcolor=lightgray](0,0){3}{98}{251}
         \pswedge[fillstyle=solid,fillcolor=brown](0,0){3}{251}{360}      
      \end{pspicture}
      \begin{pspicture}(-3.5,-3)(2.5,3.5)
         \pscircle(0,0){3}
         \pswedge[fillstyle=solid,fillcolor=Gold](0,0){3}{0}{80}
         \pswedge[fillstyle=solid,fillcolor=lightgray](0,0){3}{80}{170}
         \pswedge[fillstyle=solid,fillcolor=brown](0,0){3}{170}{360}      
      \end{pspicture}} \\
      \quad {\blue Cyclisme \hfill Tir \hfill Canoé-kayak} \hspace*{4mm} \\ \medskip
      On remarque par exemple que chacun de ces sports à une couleur très dominante : l'or pour le cyclisme, l'argent pour le tir et le bronze pour le canoé-kayak.  
\end{corrige} 

\end{colonne*exercice}


%%%%%%%%%%%%%%%%%%%%%%%%%%%%%%%%%%%%
%%%%%%%%%%%%%%%%%%%%%%%%%%%%%%%%%%%%
\Recreation

\begin{enigme}[Le tableur]
   Un tableur est un logiciel d'édition et de présentation de tableaux. Il comporte des \textbf{feuilles de calcul} composées de multiples lignes et colonnes formant des \textbf{cellules}. Chaque cellule est repérée par son adresse : une lettre désignant la colonne et un numéro désignant la ligne. Par exemple, la cellule {\bf A1} fait référence à la colonne A ligne numéro 1. \medskip

\partie[écrire dans une cellule] 
   Ouvrir une nouvelle feuille de tableur, écrire les textes suivants et appuyer sur entrée. 
   $$\text{Dans la cellule A1 : }  \Cell{\texttt{1+2}} \qquad ; \qquad \text{ Dans la cellule A2 :} \Cell{\texttt{=1+2}}$$
   Quelle est la différence entre ces deux écritures ? Quelle est la différence d'interprétation du tableur ? \par
   \Lignespointilles{2} \medskip

\partie[utilisation du tableaur-grapheur]
   Le ministère de la culture a effectué une enquête sur les équipements utilisés le plus souvent pour se connecter à Internet en 2019, en fonction de l'âge des utilisateurs.
   \begin{center}
      \begin{Tableur}[Bandeau=false,Colonnes=7,LargeurUn=5.5,Largeur=4.8]
         & 12-17 ans & 18-24 ans & 25-39 ans	 & 40-59 ans & 60-69 ans & 70 ans + \\
         Smartphone & 78\,\% & 89\,\% & 79\,\% & 45\,\% & 24\,\% & 12\,\% \\
         Ordinateur & 16\,\% & 10\,\% & 15\,\% & 42\,\% & 45\,\% & 38\,\% \\
      \end{Tableur}
   \end{center}
   On souhaite représenter, sous différentes formes, les résultats de cette enquête.
   \begin{enumerate}
      \item Reproduire ce tableau dans un tableur. \\
         {\it Enregistre ton travail dans un fichier nommé \texttt{Nom\_prenom\_classe\_tableur\_grapheur}.}
      \item Sélectionner le tableau en entier (cellules A1 à G3) et insérer trois graphiques : un diagramme en barres verticales (titre : graphique 1), un diagramme en barres horizontales (graphique 2) et un diagramme circulaire (graphique 3).
      \item Laquelle de ces trois représentations vous parait-elle la plus lisible, pourquoi ? \Lignespointilles{2}
      \item Choisir deux autres représentations et les décrire brièvement. \par
         Graphique 4 : \Lignespointilles{3} \par
         Graphique 5 : \Lignespointilles{3} \par
      \item Créer un graphique 6 de votre choix représentant uniquement les données concernant les ordinateurs.
      \item Répondre aux questions suivantes :
         \begin{enumerate}
            \item Quelle tranche d'âge utilise le plus son smartphone pour aller sur Internet ? \pointilles
            \item Selon toi, pourquoi \og les jeunes \fg{} utilisent-ils plus un smartphone qu'un ordinateur pour aller sur Internet ? \par
               \Lignespointilles{3}
            \item Selon toi, pourquoi \og les seniors \fg{} utilisent-ils plus un ordinateur qu'un smartphone pour aller sur Internet ? \par
               \Lignespointilles{3}         
         \end{enumerate}
   \end{enumerate}
\end{enigme}

